\section{引言：推理时引导生成的动机}

在前几周的课程中，我们讨论了扩散模型（Diffusion Model）和流模型（Flow Model）的基本训练和生成过程，以及如何通过推理时技术或微调来加速生成。本节课的核心问题转向另一个方向：\textbf{如何利用预训练模型，在推理时注入引导信号，使生成结果满足特定的用户偏好或约束条件？}

\begin{importantbox}{本讲核心问题}
给定一个预训练的（无条件）扩散/流模型，如何在不重新训练的前提下，在推理时引导生成过程，使输出的数据点满足特定条件——例如与给定的部分观测一致、服从某种布局约束或风格约束？
\end{importantbox}

当前存在大量预训练的扩散模型，涵盖图像生成、视频生成、分子设计、机器人运动规划等多种模态。这些模型通常需要海量数据和算力来训练，普通研究者难以从头训练。因此，一个自然的问题是：\textbf{能否仅利用预训练模型做更多事情？}

扩散/流模型的一个关键优势在于其\textbf{迭代生成}的特性。与 GAN 和 VAE 等一步生成模型不同，扩散模型在从噪声 $x_T$ 到干净数据 $x_0$ 的轨迹中，每一步都可以被"操控"：我们可以修改预测的速度或噪声，使轨迹弯曲到我们期望的终点。

\begin{knowledgebox}{回顾：Classifier Guidance 与 Classifier-Free Guidance}
在之前的课程中，我们已经见过两种在生成过程中注入引导的方法：
\begin{itemize}
    \item \textbf{Classifier Guidance}：训练一个额外的分类器网络，利用其对中间噪声数据 $x_t$ 的梯度来提供似然分数（likelihood score），从而调整生成轨迹。
    \item \textbf{Classifier-Free Guidance (CFG)}：用同一个网络同时学习条件噪声和无条件噪声，通过两者的差异隐式地计算似然分数，并用引导强度 $w$ 控制条件信号的强弱。
\end{itemize}
这两种方法都需要\textbf{训练}——要么训练额外的分类器，要么在训练噪声预测器时加入条件-无条件的配对数据。本讲讨论的是\textbf{无需额外训练}的推理时引导方法。
\end{knowledgebox}

\subsection{推理时引导的三大类别}

讲者将推理时引导技术分为三大类：

\begin{enumerate}
    \item \textbf{Score Manipulation（分数操控）}：直接对去噪过程中的分数函数进行实用性修改，无严格理论保证，但在实践中效果良好。
    \item \textbf{基于逆问题框架的引导}：将条件生成建模为逆问题（Inverse Problem），利用贝叶斯公式对似然分数进行近似计算，具有一定的理论基础。
    \item \textbf{其他高级技术}：将在下一讲（Lecture 12）讨论。
\end{enumerate}

本讲主要覆盖前两类方法。

\subsection{本章小结}

推理时引导生成是扩散模型的一个核心优势，源于其迭代生成的特性。本讲将从实用的分数操控技术出发，逐步深入到具有理论支撑的逆问题求解框架。


